% Part: first-order-logic
% Chapter: model-theory
% Section: partial-iso

\documentclass[../../../include/open-logic-section]{subfiles}

\begin{document}

\olfileid[pt]{mod}{bas}{dlo}
\section{Ordens Lineares Densas}

\begin{defn}
  Uma \emph{ordem linear densa sem extremos} é !!a{structure}
  $\Struct{M}$ para a !!{language} que contém um único
  !!{predicate} binário~$<$ e satisfaz as seguintes sentenças:
  \begin{enumerate}
  \item $\lforall[x][\lnot x < x]$;
  \item $\lforall[x][\lforall[y][\lforall[z][(x < y \lif (y < z \lif x
    <z ))]]]$;
  \item $\lforall[x][\lforall[y][(x< y \lor \eq[x][y] \lor y < x)]]$;
  \item $\lforall[x][\lexists[y][x < y]]$;
  \item $\lforall[x][\lexists[y][y < x]]$;
  \item $\lforall[x][\lforall[y][(x < y \lif \lexists[z][(x < z \land
        z < y))]]]$.
 \end{enumerate}
\end{defn}

\begin{thm}\ollabel{thm:cantorQ}
  Quaisquer duas ordens lineares densas !!{enumerable} sem
  extremos são isomorfas.
\end{thm}

\begin{proof}
  Sejam $\Struct{M_1}$ e $\Struct{M_2}$ ordens lineares densas !!{enumerable}
  sem extremos, com ${<_1} = \Assign{<}{M_1}$ e ${<_2} =
  \Assign{<}{M_2}$, e seja $\PIso{I}$ o conjunto de todos os isomorfismos
  parciais entre elas. $\PIso{I}$ não é vazio, pois ao menos
  $\emptyset \in \PIso{I}$. Mostraremos que $\PIso{I}$ satisfaz a
  propriedade de ida e volta. Então, $\Struct{M_1} \iso[p] \Struct{M_2}$,
  e o teorema decorre de \olref[pis]{thm:p-isom1}.

  Para mostrar que $\PIso{I}$ satisfaz a propriedade de ida, seja $p \in
  \PIso{I}$, e suponha que $p(a_i) = b_i$ para $i = 1$, \dots,~$n$ e,
  sem perda de generalidade, que $a_1 <_1 a_2 <_1 \cdots <_1
  a_n$. Dado $a \in \Domain{M_1}$, encontre $b \in \Domain{M_2}$ da
  seguinte maneira:
  \begin{enumerate}
  \item se $a <_1 a_1$, tome $b \in \Domain{M_2}$ tal que $b <_2
    b_1$;
  \item se $a_n <_1 a$, tome $b \in \Domain{M_2}$ tal que $b_n <_2 b$;
 \item se $a_i <_1 a <_1 a_{i+1}$ para algum $i$, tome $b \in
   \Domain{M_2}$ tal que $b_i <_2 b <_2 b_{i+1}$.
  \end{enumerate}
  Sempre é possível encontrar um $b$ com a propriedade desejada, pois
  $\Struct{M_2}$ é uma ordem linear densa sem extremos. Defina
  $q = p \cup \{ \langle a, b \rangle \}$ de modo que $q \in \PIso{I}$ seja
  a extensão desejada de $p$. Isso estabelece a propriedade de ida.
  A propriedade de volta é análoga. Portanto, $\Struct{M_1} \iso[p]
  \Struct{M_2}$; por \olref[pis]{thm:p-isom1}, $\Struct{M_1} \iso
  \Struct{M_2}$.
\end{proof}

\begin{prob}
  Complete a demonstração de \olref[mod][bas][dlo]{thm:cantorQ}
  verificando que $\PIso{I}$ satisfaz a propriedade de volta.
\end{prob}

\begin{rem}
  Seja $\Struct{S}$ qualquer ordem linear densa !!{enumerable} sem
  extremos. Então (por \olref{thm:cantorQ}) $\Struct{S} \iso
  \Struct{Q}$, em que $\Struct{Q} = (\Rat, <)$ é a ordem linear densa !!{enumerable}
  cujo domínio é o conjunto $\Rat$ dos números racionais.
  Agora, considere novamente a !!{structure}~$\Struct{R} =
  (\Real, <)$ de \olref[thm]{remark:R}. Vimos que existe
  !!a{enumerable} !!{structure}~$\Struct{S}$ tal que $\Struct{R}
  \elemequiv \Struct{S}$. Mas $\Struct{S}$ é !!a{enumerable} ordem
  linear densa sem extremos e, portanto, é isomorfa (e,
  consequentemente, elementarmente equivalente) à !!{structure}~$\Struct{Q}$. Pela
  transitividade da equivalência elementar, $\Struct{R} \elemequiv
  \Struct{Q}$. (Poderíamos ter demonstrado isso diretamente estabelecendo
  $\Struct{R} \iso[p] \Struct{Q}$ pelo mesmo argumento de ida e volta.)
\end{rem}
\end{document}
